%!TEX root = ../Thesis.tex

% This information is used in titlepage, colophon, preface and hyperref setup (pdf metainfo), and other options.

%\def\thesistypeabbr{B.Eng.}
%\def\thesistype    {Bachelor of Engineering}
%\def\thesistypeabbr{B.Sc.Eng.}
%\def\thesistype    {Bachelor of Science in Engineering}
\def\thesistypeabbr{M.Sc.}
\def\thesistype    {Master of Science in Engineering}
%\def\thesistypeabbr{Ph.D.}
%\def\thesistype    {Doctor of Philosophy}

\def\thesisauthor  {John Doe (sABCDEF)}
\def\thesistitle   {Laursen's XeLaTeX thesis template}
\def\thesissubtitle{A \emph{nice} thesis template}
\def\thesislocation{Kongens Lyngby}

\def\papersize    {b5paper} % Final papersize (b5paper/a4paper), recommended papersize for DTU Compute is b5paper
\def\showtrims    {false} % Print on larger paper than \papersize and show trim marks (true/false)?

\def\showtodos    {true}  % Show todos (true/false)?
\def\confidential {false} % Confidential thesis (true/false)?



%!TEX TS-program = xelatex
%!TEX root = ../Thesis.tex
\RequirePackage[l2tabu,orthodox]{nag} % Old habits die hard

\newcommand{\papersizeswitch}[3]{\ifnum\strcmp{\papersize}{#1}=0#2\else#3\fi}
\papersizeswitch{b5paper}{\def\classfontsize{10pt}}{\def\classfontsize{12pt}}

\documentclass[\classfontsize,\papersize,twoside,showtrims,extrafontsizes]{memoir}

% %!TEX root = ../Thesis.tex
% \RequirePackage[l2tabu,orthodox]{nag} % Old habits die hard

% \newcommand{\papersizeswitch}[3]{\ifnum\strcmp{\papersize}{#1}=0#2\else#3\fi}
% \papersizeswitch{b5paper}{\def\classfontsize{10pt}}{\def\classfontsize{12pt}}

% \documentclass[\classfontsize,\papersize,twoside,showtrims,extrafontsizes]{memoir}
\RequireXeTeX

\showtrimsoff
\papersizeswitch{b5paper}{
    % Stock and paper layout
    \pagebv
    \setlrmarginsandblock{26mm}{20mm}{*}
    \setulmarginsandblock{35mm}{30mm}{*}
    \setheadfoot{8mm}{10mm}
    \setlength{\headsep}{7mm}
    \setlength{\marginparwidth}{18mm}
    \setlength{\marginparsep}{2mm}
}{
    \papersizeswitch{a4paper}{
        \pageaiv
        \setlength{\trimtop}{0pt}
        \setlength{\trimedge}{\stockwidth}
        \addtolength{\trimedge}{-\paperwidth}
        \settypeblocksize{634pt}{448.13pt}{*}
        \setulmargins{4cm}{*}{*}
        \setlrmargins{*}{*}{0.66}
        \setmarginnotes{17pt}{51pt}{\onelineskip}
        \setheadfoot{\onelineskip}{2\onelineskip}
        \setheaderspaces{*}{2\onelineskip}{*}
    }{
    }
}
\ifnum\strcmp{\showtrims}{true}=0
    % For printing B5 on A4 with trimmarks
    \showtrimson
    \papersizeswitch{b5paper}{\stockaiv}{\stockaiii}
    \setlength{\trimtop}{\stockheight}
    \addtolength{\trimtop}{-\paperheight}
    \setlength{\trimtop}{0.5\trimtop}
    \setlength{\trimedge}{\stockwidth}
    \addtolength{\trimedge}{-\paperwidth}
    \setlength{\trimedge}{0.5\trimedge}
    
    % bigger todos if trim marks
    \setmarginnotes{10pt}{95pt}{\onelineskip}

    \trimLmarks
    
    % put jobname in left top trim mark
    \renewcommand*{\tmarktl}{%
      \begin{picture}(0,0)
        \unitlength 1mm
        \thinlines
        \put(-2,0){\line(-1,0){18}}
        \put(0,2){\line(0,1){18}}
        \put(3,15){\normalfont\ttfamily\fontsize{8bp}{10bp}\selectfont\jobname\ \
          \today\ \ 
          \printtime\ \ 
          Page \thepage}
      \end{picture}}

    % Remove middle trim marks for cleaner layout
    \renewcommand*{\tmarktm}{}
    \renewcommand*{\tmarkml}{}
    \renewcommand*{\tmarkmr}{}
    \renewcommand*{\tmarkbm}{}
\fi

\checkandfixthelayout                 % Check if errors in paper format!
\sideparmargin{outer}                 % Put sidemargins in outer position (why the fuck is this option not default by the class?)

% Large environments
\usepackage{microtype}
\usepackage{mathtools}
\usepackage{listings}                 % Source code printer for LaTeX
\usepackage{tikz}

% Links
\usepackage[hyphens]{url}             % Allow hyphens in URL's
\usepackage[unicode=false,psdextra]{hyperref}                 % References package

% Graphics and colors
\usepackage{graphicx}                 % Including graphics and using colours
\usepackage{xcolor}                   % Defined more color names
\usepackage{eso-pic}                  % Watermark and other bag
\usepackage{preamble/dtucolors}
\graphicspath{{graphics/}}

% Language
\usepackage{polyglossia}    % multilingual typesetting and appropriate hyphenation
\setdefaultlanguage{english}
\usepackage{csquotes}       % language sensitive quotation facilities

% Bibliography (references)
\usepackage[backend=biber,
            style=alphabetic,
            %backref=true,
            abbreviate=false,
            dateabbrev=false,
            alldates=long]{biblatex}

% Floating objets, captions and references
\usepackage{flafter}  % floats is positioned after or where it is defined! 
%\setfloatlocations{figure}{bhtp}   % Set floats for all figures
%\setfloatlocations{table}{bhtp}    % Set floats for all tables
%\setFloatBlockFor{section}         % Typeset floats before each section
\usepackage[noabbrev,nameinlink,capitalise]{cleveref} % Clever references. Options: "fig. !1!" --> "!Figure 1!"
\hangcaption
\captionnamefont{\bfseries}
\subcaptionlabelfont{\bfseries}
\newsubfloat{figure}
\newsubfloat{table}
%\letcountercounter{figure}{table}         % Consecutive table and figure numbering
%\letcountercounter{lstlisting}{table}     % Consecutive table and listings numbering
\captiontitlefinal{.}
% strip things from equation references, making them "(1)" instead of "Equation~1"
% from http://tex.stackexchange.com/questions/122174/how-to-strip-eq-from-cleveref
\crefformat{equation}{(#2#1#3)}
\crefrangeformat{equation}{(#3#1#4) to~(#5#2#6)}
\crefmultiformat{equation}{(#2#1#3)}%
{ and~(#2#1#3)}{, (#2#1#3)}{ and~(#2#1#3)}

% Table of contents (TOC)
\setcounter{tocdepth}{1}              % Depth of table of content
\setcounter{secnumdepth}{2}           % Depth of section numbering
\setcounter{maxsecnumdepth}{3}        % Max depth of section numbering

% Todos
\usepackage{totcount}                 % For total counting of counters
\def\todoshowing{}
\ifnum\strcmp{\showtodos}{false}=0
    \def\todoshowing{disable}
\fi
\usepackage[colorinlistoftodos,\todoshowing]{todonotes} % Todonotes package for nice todos
\newtotcounter{todocounter}           % Creates counter in todo
\let\oldtodo\todo
\newcommand*{\newtodo}[2][]{\stepcounter{todocounter}\oldtodo[#1]{\thesection~(\thetodocounter)~#2}}
\let\todo\newtodo
\let\oldmissingfigure\missingfigure
\newcommand*{\newmissingfigure}[2][]{\stepcounter{todocounter}\oldmissingfigure[#1]{\thesection~(\thetodocounter)~#2}}
\let\missingfigure\newmissingfigure
\makeatletter
\newcommand*{\mylistoftodos}{% Only show list if there are todos
\if@todonotes@disabled
\else
    \ifnum\totvalue{todocounter}>0
        \markboth{\@todonotes@todolistname}{\@todonotes@todolistname}
        \phantomsection\todototoc
        \listoftodos
    \else
    \fi
\fi
}
\makeatother
\newcommand{\lesstodo}[2][]{\todo[color=green!40,#1]{#2}}
\newcommand{\moretodo}[2][]{\todo[color=red!40,#1]{#2}}

% Chapterstyle
\makeatletter
\makechapterstyle{mychapterstyle}{
    \chapterstyle{default}
    \def\format{\normalfont\sffamily}

    \setlength\beforechapskip{0mm}

    \renewcommand*{\chapnamefont}{\format\HUGE}
    \renewcommand*{\chapnumfont}{\format\fontsize{54}{54}\selectfont}
    \renewcommand*{\chaptitlefont}{\format\fontsize{42}{42}\selectfont}

    \renewcommand*{\printchaptername}{\chapnamefont\MakeUppercase{\@chapapp}}
    \patchcommand{\printchaptername}{\begingroup\color{dtugray}}{\endgroup}
    \renewcommand*{\chapternamenum}{\space\space}
    \patchcommand{\printchapternum}{\begingroup\color{dtured}}{\endgroup}
    \renewcommand*{\printchapternonum}{%
        \vphantom{\printchaptername\chapternamenum\chapnumfont 1}
        \afterchapternum
    }

    \setlength\midchapskip{1ex}

    \renewcommand*{\printchaptertitle}[1]{\raggedleft \chaptitlefont ##1}
    \renewcommand*{\afterchaptertitle}{\vskip0.5\onelineskip \hrule \vskip1.3\onelineskip}
}
\makeatother
\chapterstyle{mychapterstyle}

% Header and footer
\def\hffont{\sffamily\small}
\makepagestyle{myruled}
\makeheadrule{myruled}{\textwidth}{\normalrulethickness}
\makeevenhead{myruled}{\hffont\thepage}{}{\hffont\leftmark}
\makeoddhead{myruled}{\hffont\rightmark}{}{\hffont\thepage}
\makeevenfoot{myruled}{}{}{}
\makeoddfoot{myruled}{}{}{}
\makepsmarks{myruled}{
    \nouppercaseheads
    \createmark{chapter}{both}{shownumber}{}{\space}
    \createmark{section}{right}{shownumber}{}{\space}
    \createplainmark{toc}{both}{\contentsname}
    \createplainmark{lof}{both}{\listfigurename}
    \createplainmark{lot}{both}{\listtablename}
    \createplainmark{bib}{both}{\bibname}
    \createplainmark{index}{both}{\indexname}
    \createplainmark{glossary}{both}{\glossaryname}
}
\pagestyle{myruled}
\copypagestyle{cleared}{myruled}      % When \cleardoublepage, use myruled instead of empty
\makeevenhead{cleared}{\hffont\thepage}{}{} % Remove leftmark on cleared pages

\makeevenfoot{plain}{}{}{}            % No page number on plain even pages (chapter begin)
\makeoddfoot{plain}{}{}{}             % No page number on plain odd pages (chapter begin)

% \*section, \*paragraph font styles
\setsecheadstyle              {\huge\sffamily\raggedright}
\setsubsecheadstyle           {\LARGE\sffamily\raggedright}
\setsubsubsecheadstyle        {\Large\sffamily\raggedright}
%\setparaheadstyle             {\normalsize\sffamily\itseries\raggedright}
%\setsubparaheadstyle          {\normalsize\sffamily\raggedright}


% Hypersetup
\hypersetup{
    pdfauthor={\thesisauthor{}},
    pdftitle={\thesistitle{}},
    pdfsubject={\thesissubtitle{}},
    pdfdisplaydoctitle,
    bookmarksnumbered=true,
    bookmarksopen,
    breaklinks,
    linktoc=all,
    plainpages=false,
    unicode=true,
    colorlinks=false,
    citebordercolor=dtured,           % color of links to bibliography
    filebordercolor=dtured,           % color of file links
    linkbordercolor=dtured,           % color of internal links (change box color with linkbordercolor)
    urlbordercolor=s13,               % color of external links
    hidelinks,                        % Do not show boxes or colored links.
}
% Hack to make right pdfbookmark link. The normal behavior links just below the chapter title. This hack put the link just above the chapter like any other normal use of \chapter.
% Another solution can be found in http://tex.stackexchange.com/questions/59359/certain-hyperlinks-memoirhyperref-placed-too-low
\makeatletter
\renewcommand{\@memb@bchap}{%
  \ifnobibintoc\else
    \phantomsection
    \addcontentsline{toc}{chapter}{\bibname}%
  \fi
  \chapter*{\bibname}%
  \bibmark
  \prebibhook
}
\let\oldtableofcontents\tableofcontents
\newcommand{\newtableofcontents}{
    \@ifstar{\oldtableofcontents*}{
        \phantomsection\addcontentsline{toc}{chapter}{\contentsname}\oldtableofcontents*}}
\let\tableofcontents\newtableofcontents
\makeatother

% Confidential
\newcommand{\confidentialbox}[1]{
    \put(0,0){\parbox[b][\paperheight]{\paperwidth}{
        \begin{vplace}
            \centering
            \scalebox{1.3}{
                \begin{tikzpicture}
                    \node[very thick,draw=red!#1,color=red!#1,
                          rounded corners=2pt,inner sep=8pt,rotate=-20]
                          {\sffamily \HUGE \MakeUppercase{Confidential}};
                \end{tikzpicture}
            }
        \end{vplace}
    }}
}

% Prefrontmatter
\newcommand{\prefrontmatter}{
    \pagenumbering{alph}
    \ifnum\strcmp{\confidential}{true}=0
        \AddToShipoutPictureBG{\confidentialbox{10}}   % 10% classified box in background on each page
        \AddToShipoutPictureFG*{\confidentialbox{100}} % 100% classified box in foreground on first page
    \fi
}

% DTU frieze
\newcommand{\frieze}{%
    \AddToShipoutPicture*{
        \put(0,0){
            \parbox[b][\paperheight]{\paperwidth}{%
                \includegraphics[trim=130mm 0 0 0,width=0.9\textwidth]{DTU-frise-SH-15}
                \vspace*{2.5cm}
            }
        }
    }
}

% This is a double sided book. If there is a last empty page lets use it for some fun e.g. the frieze.
% NB: For a fully functional hack the \clearpage used in \include does some odd thinks with the sequence numbering. Thefore use \input instead of \include at the end of the book. If bibliography is used at last everything should be ok.
\makeatletter
% Adjust so gatherings is allowd for single sheets too! (hacking functions in memoir.dtx)
\patchcmd{\leavespergathering}{\ifnum\@memcnta<\tw@}{\ifnum\@memcnta<\@ne}{
    \leavespergathering{1}
    % Insert the frieze
    \patchcmd{\@memensuresigpages}{\repeat}{\repeat\frieze}{}{}
}{}
\makeatother

%!TEX root = ../Thesis.tex

% Text fonts (http://www.macfreek.nl/memory/Fonts_in_LaTeX)
% Install fonts from /usr/local/texlive/<version>/texmf-dist/fonts/opentype/public
\usepackage{fontspec}

% Sans-serif font
\setsansfont[
    Ligatures=TeX,
    Extension=.otf,
    UprightFont=*-regular,
    BoldFont=*-bold,
    ItalicFont=*-italic,
    BoldItalicFont=*-bolditalic,
    %SlantedFont=,
    %BoldSlantedFont=,
    %SmallCapsFont=
    Scale=0.8      % Adjustmens when using math in sections
]{texgyreadventor}
%\setsansfont[Ligatures=TeX]{Neo Sans Intel}    % Neo Sans Intel – Like DTU font but more symbols
%\setsansfont[
%    Ligatures=TeX,
%    Scale=0.8
%]{NeoSans}           % NeoSans – DTU font (missing `+' symbols and other)
%\setsansfont[Ligatures=TeX]{CMU Sans Serif}    % Computer Modern Unicode font
%\setsansfont[Ligatures=TeX]{Latin Modern Sans} % Latin Modern Sans serif font

% Use this for more convienent sans serif font in math mode.
%\setmathsf{Latin Modern Sans}


%!TEX root = ../Thesis.tex

% Content specific packages.

\usepackage{blindtext}
\usepackage{algorithm}
\usepackage{algpseudocode}
%\usepackage{pgfplots}                 % Plot tools
%\usetikzlibrary{
%    arrows,
%    matrix,
%    positioning,
%    shapes,
%    topaths,
%}
%\pgfplotsset{compat=1.7}

% Listings
\lstset{
    basicstyle=\footnotesize\ttfamily,% the size of the fonts that are used for the code
    breakatwhitespace=false,          % sets if automatic breaks should only happen at whitespace
    breaklines=true,                  % sets automatic line breaking
    captionpos=b,                     % sets the caption-position to bottom
    commentstyle=\color{s14a},        % comment style
    deletekeywords={},                % if you want to delete keywords from the given language
    escapeinside={\%*}{*)},           % if you want to add LaTeX within your code
    frame=single,                     % adds a frame around the code
    keywordstyle=\bfseries\ttfamily\color{s09}, % keyword style
    language=Python,                  % the language of the code
    morekeywords={*,...},             % if you want to add more keywords to the set
    numbers=left,                     % where to put the line-numbers; possible values are (none, left, right)
    numbersep=5pt,                    % how far the line-numbers are from the code
    numberstyle=\sffamily\tiny\color{dtugray}, % the style that is used for the line-numbers
    rulecolor=\color{dtugray},        % if not set, the frame-color may be changed on line-breaks within not-black text (e.g. comments (green here))
    showspaces=false,                 % show spaces everywhere adding particular underscores; it overrides 'showstringspaces'
    showstringspaces=false,           % underline spaces within strings only
    showtabs=false,                   % show tabs within strings adding particular underscores
    stepnumber=1,                     % the step between two line-numbers. If it's 1, each line will be numbered
    stringstyle=\color{s07},          % string literal style
    tabsize=2,                        % sets default tabsize to 2 spaces
    title=\lstname,                   % show the filename of files included with \lstinputlisting; also try caption instead of title
}



\addbibresource{bibliography/Bibliography.bib}

\begin{document}

\prefrontmatter
%!TEX root = ../Thesis.tex 
\thispagestyle{empty}             % No page numbers
\calccentering{\unitlength}
\begin{adjustwidth*}{\unitlength}{-\unitlength}
    \begin{adjustwidth}{-0.5cm}{-0.5cm}
        \sffamily
        \begin{flushright}
            \thesistypeabbr{} Thesis\\*[0cm]
            \thesistype{}\\
        \end{flushright}
        \vspace*{\fill}
        \noindent
        \includegraphics[width=0.55\textwidth]{DTU-EE}\\*[0.5cm]
        \HUGE \thesistitle{}\\*[0.2cm]
        \Huge \thesissubtitle{}\\*[0.8cm]
        \parbox[b]{0.5\linewidth}{%
            \LARGE 
            \thesisauthor{}\\*[1.2cm]
            \Large
            \thesislocation{} \the\year
        }
        \hfill\includegraphics[scale=0.7]{DTU-logo-CMYK}
    \end{adjustwidth}
\end{adjustwidth*}
\normalfont
\normalsize

\cleartoevenpage
%!TEX root = ../Thesis.tex
\thispagestyle{empty} % No page numbers
\frieze
\vspace*{\fill}
\noindent
\sffamily
\small
\textbf{DTU Compute}\\
\textbf{Department of Applied Mathematics and Computer Science}\\
\textbf{Technical University of Denmark}\\
\\
Matematiktorvet\\
Building 303B\\
2800 Kongens Lyngby, Denmark\\
Phone +45 4525 3031\\
compute@compute.dtu.dk\\
www.compute.dtu.dk\\
\normalsize
\normalfont
\vspace*{2.5cm}

\clearforchapter

\frontmatter
%!TEX root = ../Thesis.tex


% Summary and foreword
% Most readers will turn first to the summary (or abstract). Use it as an opportunity to spur the reader’s interest. The summary should highlight the main points from your work, especially the thesis statement, methods (if applicable), findings and conclusion. However, the summary does not need to cover every aspect of your work. The main objective is to give the reader a good idea of what the thesis is about.

% The summary should be completed towards the end; when you are able to overview your project as a whole. It is nevertheless a good idea to work on a draft continuously. Writing a good summary can be difficult, since it should only include the most important points of your work. But this is also why working on your summary can be so useful – it forces you to identify the key elements of your writing project.

% There are usually no formal requirements for forewords, but it is common practice to thank your supervisors, informants, and others who have helped and supported you. If you have received any grants or research residencies, you should also acknowledge these.

% Note: Shorter assignments do not require abstracts and forewords. 


\chapter{Abstract}

In the following thesis, the author provides a comprehensive description of a process, decision making and final implementation of a master control unit for the electric racing car. The workload covered is a patronate over design of intelligent sensors/controllers, implementation of event-based CAN library for real-time system, control motor controllers over the CANOpen protocol, communication with a battery management system, logging system state and events, implementing safe and reliable boot/shutdown sequence and finally on the top of it controlling the car under normal operation including adaptive cooling pumps control, torque vectoring as replacement for mechanical differential and regenerative braking.
%!TEX root = ../Thesis.tex
\chapter{Preface}
This thesis was prepared at the department of Electrical Engineering at the Technical University of Denmark in fulfilment of the requirements for acquiring a master degree in Electrical Engineering.

\vfill

{
\centering
    \thesislocation{}, \today\\[1cm]
    \hspace{3cm}\includegraphics[scale=0.4]{Signature}\\[1cm]
\begin{flushright}
    \thesisauthor{}
\end{flushright}
}
%!TEX root = ../Thesis.tex
\chapter{Acknowledgements}

I would like to thank my DTU's supervisor Nenad Mijatovic for giving me a lot of trust and understanding, so I could organise my working time in the most suitable way. There were many perturbations on the way and although stress and unlucky events, I have never been criticised or nag but rather admired for my work. This gave me the power to go through the moments where I doubted if my skills and expertise are up to the task.

In the final stage of work, I needed to move back to Korea, where I am writing these words right now. The need was not only dictated by KAIST graduation requirements but I believed that it would help me in finding motivation for writing. All my hopes were overfilled by Profesor Naehyuck Chang (KAIST's supervisor) who provided me with a great working environment and have been the first to review my this thesis. His opinion was valuable for me, he gave me hope and pointed further direction.

Of course, I would like to thank the whole team working on the car, I believe it was a great experience for us all. Here I would like to give special acknowledgement to Martin van der Bijl who kept me company during never-ending days in the laboratory, it was really nice to have someone to share my thoughts with time to time.

Moreover, I would like to thank my boss (to cover my expenses of living in aboard I needed to simultaneously work through the whole period of my master studies). I have been offered position when due to financial reasons I been considered dropping studies. I felt anxious as it was my first long-term job, however, I was given a lot of recognition and flexibility which made it possible for me to attend and complete this degree.

Last but not least I would like to thank my family for their never-ending mental support and never questioning my decisions. You are the root of each and every one of my little success.  
\clearforchapter
\tableofcontents
\clearforchapter
\mylistoftodos

\mainmatter

% 1.2 Defining the scope of your thesis
% One of the first tasks of a researcher is defining the scope of a study, i.e., its area (theme, field) and the amount of information to be included. Narrowing the scope of your thesis can be time-consuming. Paradoxically, the more you limit the scope, the more interesting it becomes. This is because a narrower scope lets you clarify the problem and study it at greater depth, whereas very broad research questions only allow a superficial treatment.

% The research question can be formulated as one main question with (a few) more specific sub-questions or in the form of a hypothesis that will be tested.

% Your research question will be your guide as your writing proceeds. If you are working independently, you are also free to modify it as you go along.

% How do you know that you have drafted a research question? Most importantly, a research question is something that can be answered. If not, you have probably come up with a theme or field, not a question.

% Some tips:

% Use interrogative words: how, why, which (factors/situations) etc.
% Some questions are closed and only invoke concrete/limited answers. Others will open up for discussions and different interpretations.
% Asking “What …?” is a more closed question than asking “How?” or “In what way?”
% Asking “Why” means you are investigating what causes of a phenomenon. Studying causality is methodologically demanding.
% Feel free to pose partially open questions that allow discussions of the overall theme, e.g., “In what way …?”; “How can we understand [a particular phenomenon]?”
% Try to condense your research question into one general question – and perhaps a few more specific sub-questions (two or three will usually suffice).


% 2. Theory section
% The theory used in an empirical study is meant to shed light on the data in a scholarly or scientific manner. It should give insights not achievable by ordinary, everyday reflections. The main purpose of using theory is to analyse and interpret your data. Therefore, you should not present theoretical perspectives that are not being put to use. Doing so will create false expectations, and suggests that your work is incomplete.

% Not all theses have a separate theory section. In the IMRaD format the theory section is included in the introduction, and the second chapter covers the methods used.

% What kind of theory should you choose? Since the theory is the foundation for your data analysis it can be useful to select a theory that lets you distinguish between, and categorise different phenomena. Other theories let you develop the various nuances of a phenomenon. In other words, you have a choice of either reducing the complexity of your data or expanding upon something that initially looks simple.

% How much time and space should you devote to the theory chapter? This is a difficult question. Some theses dwell too long on theory and never get to the main point: the analysis and discussion. But it is also important to have read enough theory to know what to look for when collecting data. The nature of your research should decide: Some studies do not require much theory, but put more emphasis on the method, while other studies need a rich theory section to enable an interesting discussion.

% 3. Method section
% In a scholarly research article, the section dealing with method is very important. The same applies to an empirical thesis. For students, this can be a difficult section to write, especially since its purpose may not always be clear.

% The method chapter should not iterate the contents of methodology handbooks. For example, if you have carried out interviews, you do not need to list all the different types of research interview. You also do not need to describe the differences between quantitative and qualitative methods, or list all different kinds of validity and reliability.

% What you must do is to show how your choice of design and research method is suited to answering your research question(s). Demonstrate that you have given due consideration to the validity and reliability of your chosen method. By “showing” instead of “telling”, you demonstrate that you have understood the practical meaning of these concepts. This way, the method section is not only able to tie the different parts of your thesis together, it also becomes interesting to read!

% Show the reader what you have done in your study, and explain why. How did you collect the data? Which options became available through your chosen approach?
% What were your working conditions? What considerations did you have to balance?
% Tell the reader what you did to increase the validity of your research. E.g., what can you say about the reliability in data collection? How do you know that you have actually investigated what you intended to investigate? What conclusions can be drawn on this basis? Which conclusions are certain and which are more tentative? Can your results be applied in other areas? Can you generalise? If so, why? If not, why not?
% You should aim to describe weaknesses as well as strengths. An excellent thesis distinguishes itself by defending – and at the same time criticising – the choices made.
% 4. Analysis
% Your analysis, along with your discussion, will form the high light of your thesis. In the IMRaD format, this section is titled “Results”. This is where you report your findings and present them in a systematic manner. The expectations of the reader have been built up through the other chapters, make sure you fulfill these expectations.

% To analyse means to distinguish between different types of phenomena – similar from different. Importantly, by distinguishing between different phenomena, your theory is put to work. Precisely how your analysis should appear, however, is a methodological question. Finding out how best to organise and present your findings may take some time. A good place to look for examples and inspiration is repositories for master’s theses.

% If you are analysing human actions, you may want to engage the reader’s emotions. In this case it will be important to choose analytical categories that correlate to your chosen theory. Engaging emotions is not the main point, but a way to elucidate the phenomenon so that the reader understands it in a new and better way.

% Note: Not all theses include a separate chapter for analysis.

% 5. Discussion
% In many thesis the discussion is the most important section. Make sure that you allocate enough time and space for a good discussion. This is your opportunity to show that you have understood the significance of your findings and that you are capable of applying theory in an independent manner.

% The discussion will consist of argumentation. In other words, you investigate a phenomenon from several different perspectives. To discuss means to question your findings, and to consider different interpretations. Here are a few examples of formulations that signal argumentation:

% On the one hand … and on the other …
% However …
% … it could also be argued that …
% … another possible explanation may be …
% 6. Conclusion – or summing up?

% The final section of your thesis may take one of several different forms. Some theses need a conclusion, while for others a summing up will be appropriate. The decisive factor will be the nature of your thesis statement and/or research question.

% Open research questions cannot always be answered, but if a definite answer is possible, you must provide a conclusion. The conclusion should answer your research question(s). Remember that a negative conclusion is also valid.

% A summing up should repeat the most important issues raised in your thesis (particularly in the discussion), although preferably stated in a (slightly) different way. For example, you could frame the issues within a wider context.

% Placing your thesis in perspective

% In the final section you should place your work in a wider, academic perspective and determine any unresolved questions. During the work, you may have encountered new research questions and interesting literature which could have been followed up. At this point, you may point out these possible developments, while making it clear for the reader that they were beyond the framework of your current project.

% Briefly discuss your results through a different perspective. This will allow you to see aspects that were not apparent to you at the project preparation stage
% Highlight alternative research questions that you have found in the source materials used in the project
% Show how others have placed the subject area in a wider context
% If others have drawn different conclusions from yours, this will provide you with ideas of new ways to view the research question
% Describe any unanswered aspects of your project
% Specify potential follow up and new projects
% A thesis should “bite itself in the tail”

% There should be a strong connection between your conclusion and your introduction. All the themes and issues that you raised in your introduction must be referred to again in one way or another. If you find out at this stage that your thesis has not tackled an issue that you raised in the introduction, you should go back to the introduction and delete the reference to that issue. An elegant way to structure the text is to use the same textual figure or case in the beginning as well as in the end. When the figure returns in the final section, it will have taken on a new and richer meaning through the insights you have encountered, created in the process of writing.






%!TEX root = ../Thesis.tex
\chapter{Introduction}

\begin{itemize}
    \item \textup{Upright shape}
    \item \textit{Italic shape}
    \item \textsl{Slanted shape}
    \item \textsc{Small Caps shape}
    \item \textmd{Medium series}
    \item \textbf{Bold sereies}
    \item \textrm{Roman family}
    \item \textsf{Sans serif family}
    \item \texttt{Typewriter family}
\end{itemize}

I love to write special characters like øæå indside my \TeX{} document. Also á, à, ü, û, ë, ê, î, ï could be nice. So waht about the ``π'' chracter. What about ° é ® † ¥ ü | œ π ‘ @ ö ä ¬ ∆ ‹ « © ƒ ∂ ß ª Ω … ç √ ∫ ñ µ ‚ · ¡ “ £ ∞ ™ [ ] ≠ ± '.

Some dashes - – —, and the latex form - -- ---
\begin{equation*}
    x = \mathtt{x}, \mathbf{x}, \mathit{x}, x_{1_{2_{3_{4}}}}^{1^{2^{3^{4}}}} \cdot hello * \text{hello world} ⋅ \text{my world} · \text{third world} ⊗ t
\end{equation*}

Lorem ipsum dolor sit amet, consectetur adipisicing elit, sed do eiusmod tempor incididunt ut labore et dolore magna aliqua\todo{Make a cake}. Ut enim ad minim veniam, quis nostrud exercitation ullamco laboris nisi ut aliquip ex ea commodo consequat. Duis aute irure dolor in reprehenderit in voluptate velit esse cillum dolore eu fugiat nulla pariatur. Excepteur sint occaecat cupidatat non proident, sunt in culpa qui officia deserunt mollit anim id est laborum \cite{adams1980hitchhiker}.

Mauris id quam non magna fermentum malesuada id mattis lorem. In a dapibus neque. Etiam lacus dui, malesuada ac eleifend imperdiet, imperdiet ut ipsum\todo{Do it now}. Vestibulum id ultricies est. Phasellus augue mauris, semper a luctus vel, faucibus in risus. Fusce commodo augue quis elit sagittis non viverra turpis bibendum. Nunc placerat sem non sapien malesuada malesuada ullamcorper orci luctus \cite{adams1980hitchhiker}. Morbi pharetra ligula integer mollis mi nec neque ultrices vitae volutpat leo ullamcorper. In at tellus magna. Curabitur quis posuere purus. Cum sociis natoque penatibus et magnis dis parturient montes, nascetur ridiculus mus. Suspendisse tristique placerat feugiat. Aliquam vitae est at enim auctor ultrices eleifend a urna. Donec non tincidunt felis. Maecenas at suscipit orci. See \cref{myFigure}.
\begin{figure}
    \centering
        \missingfigure[figwidth=6cm]{This is some text that is with the todo and in the figure}
    \caption[Short caption to special figure]{This is my special figure. Aliquam ultricies, arcu quis tempor rhoncus, tellus nisl tempus justo, condimentum tempor erat odio ac purus. Integer quis ipsum felis. Aliquam volutpat, leo ac consequat egestas, lectus lacus adipiscing quam, id iaculis dolor quam in erat. Phasellus tempor interdum arcu quis vestibulum}
    \label{myFigure}
\end{figure}

Fusce id suscipit sem. Aliquam venenatis nibh nec nisl luctus vel consectetur neque dapibus. Nulla feugiat egestas turpis, ac viverra eros cursus sit amet. Cras tincidunt felis vel tellus ultrices condimentum. Quisque vehicula, arcu vitae interdum dignissim, purus tortor cursus libero, sit amet accumsan quam magna in neque. Phasellus luctus leo odio. Aliquam ultricies, arcu quis tempor rhoncus, tellus nisl tempus justo, condimentum tempor erat odio ac purus. Integer quis ipsum felis. Aliquam volutpat, leo ac consequat egestas, lectus lacus adipiscing quam, id iaculis dolor quam in erat. Phasellus tempor interdum arcu quis vestibulum. Pellentesque sit amet augue purus. See \cref{myTable}.

\begin{table}
    \centering
        \begin{tabular}{c | r l}
            h & h & h \\
            e & e & e \\
        \end{tabular}
    \caption{This is a caption to the table}
    \label{myTable}
\end{table}

\section{Torquent Arcu}
Curabitur condimentum suscipit arcu, sit amet convallis urna pellentesque ac. Quisque fringilla tincidunt risus nec accumsan. Curabitur vel sagittis ante. Integer eget placerat leo. Class aptent taciti sociosqu ad litora torquent per conubia nostra, per inceptos himenaeos. Vestibulum quis risus in nulla fermentum pellentesque dictum et erat. Nulla vel pretium nunc. Integer tortor lorem, suscipit sit amet ultricies non, porta at metus. Sed pharetra, ante facilisis interdum porta, mi dolor fringilla quam, ac porttitor urna dolor quis massa. Proin viverra semper tincidunt. Vivamus pulvinar pharetra condimentum. Pellentesque rutrum mollis tellus ac scelerisque.

\begin{figure}
    \centering
        \subbottom[1 pass]{
            \missingfigure[figwidth=6cm]{This is some text that is with the todo and in the figure}
            \label{twoFigures1}
        }
        ~
        \subbottom[5 passes]{
            \missingfigure[figwidth=6cm]{This is some text that is with the todo and in the figure}
            \label{twoFigures2}
        }
        \caption{loop performance comparison}
    \label{twoFigures}
\end{figure}

\subsection{Vestibulum}
Mauris luctus sollicitudin vestibulum. Class aptent taciti sociosqu ad litora torquent per conubia nostra, per inceptos himenaeos \cref{twoFigures2}. Duis eu nisl nec turpis porttitor bibendum eget sed orci. Aliquam consequat lorem a dui viverra porta facilisis augue rutrum. Cras luctus tellus in lectus egestas eu consequat magna cursus. Aenean aliquam neque a nibh elementum ornare. Integer eleifend imperdiet commodo. Morbi auctor, dui vel laoreet congue, purus est accumsan augue, sit amet feugiat neque nisl vel lorem. Curabitur ante sem, lacinia id adipiscing quis, viverra tristique nulla. Pellentesque ullamcorper pellentesque metus varius facilisis. Cras ac dui id odio tempor scelerisque. Curabitur a egestas risus. Pellentesque quis velit in sapien accumsan auctor. Phasellus aliquam, sapien eget lobortis volutpat, libero metus porttitor nisl, sed hendrerit urna dolor nec mi. See \cref{fibonacci}.

\begin{adjustwidth*}{0cm}{-0.4cm}
\begin{lstlisting}[language=Python,caption=Fibonacci,label=fibonacci]
# This is a comment
import easy
str = "I am a string"
str2 = "Now i have an awsome string with ´ '' `` which are not TeX'ed"
str3 = "What about awsome unicode characters? Like “, π, ”, Ω, ç. \" This"
def fib(n):
    if n == 0:
        return 0
    elif n == 1:
        return 1
    else:
        return fib(n-1) + fib(n-2)
str4 = "Yes it is possible with 80 charactes. Which this string proves. Wiiii."
str5 = "It adjusts according to the spine"
\end{lstlisting}
\end{adjustwidth*}

\section{Luctus}
Praesent et pellentesque arcu. Phasellus venenatis mi eu lorem convallis et iaculis ante aliquet. Aenean rhoncus placerat metus, vel convallis leo suscipit eu. Integer dapibus venenatis commodo. Cras laoreet faucibus sem nec luctus. Class aptent taciti sociosqu ad litora torquent per conubia nostra, per inceptos himenaeos. Cras consectetur lacinia dolor at gravida. Phasellus ipsum arcu, vulputate fermentum ultricies eget, tempor eu odio. Aenean accumsan vestibulum risus a mattis. See it on \cref{modifiedminibatch}.

\begin{algorithm}
\caption{Modified mini-batch $K$-means} \label{modifiedminibatch}
\begin{algorithmic}[1]
\State Given: $K$, mini-batch size $B$, iterations $T$, dataset $X$, correlation~matrix~$\mathrm{P}$.
\State Initialize $C = \{\mathbf{c}^{(1)}, \mathbf{c}^{(2)}, \ldots, \mathbf{c}^{(K)}\}$ with random $\mathbf{x}$'es picked from $X$.
\State $A \gets B \cdot T$ sorted random indexes to $X$, denoted $a_1, a_2, \ldots, a_{B\cdot T}$.
\State $X' \gets \{\mathbf{x}^{(a_1)}, \mathbf{x}^{(a_2)}, \ldots, \mathbf{x}^{(a_{B\cdot T})}\}$ \Comment{Cache all points}
\State $\textbf{size} \gets 0$
\For {$i = 1$ to $T$}
    \State $M \gets B$ examples picked randomly from $X'$
    
    \For{$\mathbf{x} \in M$} \Comment{\textit{Assignment step}}
        \State $\textbf{d}[\textbf{x}] \gets f(C,\mathbf{x}, \mathrm{P})$ \Comment{Cache closest center}
    \EndFor
    
    \For {$\mathbf{x} \in M$} \Comment{\textit{Update step}}
        \State $\textbf{c} \gets \textbf{d[x]}$ \Comment{Get cached center for current \textbf{x}}
        \State $\textbf{size}[\textbf{c}] \gets \textbf{size}[\textbf{c}] + 1$ \Comment{Update cluster size}
        \State $\eta \gets \frac{1}{\textbf{size}[\textbf{c}]}$       \Comment{Get learning rate}
        \State $\textbf{c} \gets (1 - \eta)\textbf{c}+\eta\textbf{x}$ \Comment{Take gradient step}
    \EndFor
\EndFor
\State \Return {$C$, \textbf{size}}
\end{algorithmic}
\end{algorithm}

Fusce id suscipit sem. Aliquam venenatis nibh nec nisl luctus vel consectetur neque dapibus. Nulla feugiat egestas turpis, ac viverra eros cursus sit amet. Cras tincidunt felis vel tellus ultrices condimentum. Quisque vehicula, arcu vitae interdum dignissim, purus tortor cursus libero, sit amet accumsan quam magna in neque. Phasellus luctus leo odio. Aliquam ultricies, arcu quis tempor rhoncus, tellus nisl tempus justo, condimentum tempor erat odio ac purus. Integer quis ipsum felis. Aliquam volutpat, leo ac consequat egestas, lectus lacus adipiscing quam, id iaculis dolor quam in erat. Phasellus tempor interdum arcu quis vestibulum. Pellentesque sit amet augue purus. 
\begin{figure}[htbp]
    \centering
        \missingfigure[figwidth=6cm]{This is some text that is with the todo and in the figure}
    \caption{This is the caption I wrote}
    \label{fig:label}
\end{figure}
Curabitur condimentum suscipit arcu, sit amet convallis urna pellentesque ac. Quisque fringilla tincidunt risus nec accumsan. Curabitur vel sagittis ante. Integer eget placerat leo. Class aptent taciti sociosqu ad litora torquent per conubia nostra, per inceptos himenaeos. Vestibulum quis risus in nulla fermentum pellentesque dictum et erat. Nulla vel pretium nunc. Integer tortor lorem, suscipit sit amet ultricies non, porta at metus. Sed pharetra, ante facilisis interdum porta, mi dolor fringilla quam, ac porttitor urna dolor quis massa. Proin viverra semper tincidunt. Vivamus pulvinar pharetra condimentum. Pellentesque rutrum mollis tellus ac scelerisque.

\section{Sollicitudin vestibulum}
Mauris luctus sollicitudin vestibulum. Class aptent taciti sociosqu ad litora torquent per conubia nostra, per inceptos himenaeos. Duis eu nisl nec turpis porttitor bibendum eget sed orci. Aliquam consequat lorem a dui viverra porta facilisis augue rutrum. Cras luctus tellus in lectus egestas eu consequat magna cursus. Aenean aliquam neque a nibh elementum ornare. Integer eleifend imperdiet commodo. Morbi auctor, dui vel laoreet congue, purus est accumsan augue, sit amet feugiat neque nisl vel lorem. Curabitur ante sem, lacinia id adipiscing quis, viverra tristique nulla. Pellentesque ullamcorper pellentesque metus varius facilisis. Cras ac dui id odio tempor scelerisque. Curabitur a egestas risus. Pellentesque quis velit in sapien accumsan auctor. Phasellus aliquam, sapien eget lobortis volutpat, libero metus porttitor nisl, sed hendrerit urna dolor nec mi.

Praesent et pellentesque arcu. Phasellus venenatis mi eu lorem convallis et iaculis ante aliquet. Aenean rhoncus placerat metus, vel convallis leo suscipit eu. Integer dapibus venenatis commodo. Cras laoreet faucibus sem nec luctus. Class aptent taciti sociosqu ad litora torquent per conubia nostra, per inceptos himenaeos. Cras consectetur lacinia dolor at gravida. Phasellus ipsum arcu, vulputate fermentum ultricies eget, tempor eu odio. Aenean accumsan vestibulum risus a mattis.

\begin{adjustwidth*}{0cm}{-0.4cm}
\begin{lstlisting}[language=Python,caption=Fibonacci2,label=Fibonacci2]
# This is a comment
import easy
str = "I am a string"
str2 = "Now i have an awsome string with ´ '' `` which are not TeX'ed"
str3 = "What about awsome unicode characters? Like “, π, ”, Ω, ç. \" This"
def fib(n):
    if n == 0:
        return 0
    elif n == 1:
        return 1
    else:
        return fib(n-1) + fib(n-2)
str4 = "Yes it is possible with 80 charactes. Which this string proves. Wiiii."
str5 = "It adjusts according to the spine"
\end{lstlisting}
\end{adjustwidth*}
\include{chapters/TechnicalBackground}
\include{chapters/ExistingSolutions}
%!TEX root = ../Thesis.tex
\chapter{Overview of experimental setup}
Having not sufficient financial support enforced many compromises on design which I did not have direct influence on. The decision was made that will be running in low voltage - high current setup so we could use cost efficient lithium-ion batteries. 
Although theoretically possible cabling and connector sizing was a bit overwhelming for responsible teams.

In our final setup we used:
\begin{itemize}
    \item Two EMRAX228 motors which nominal power of about 50kW (when liquid cooled)
    \item Two emDrive500 controllers capable full usage of AM motors
    \item Arduinos with CAN shields to act as intelligent sensors/controllers
    \item Two \todo{which pump?} pumps for cooling purposes
    \item National Instruments' cRIO 9033 as a main control unit
\end{itemize}

This result in system which in simplified manner has been presented in figure \todo{figure ref}.

\todo{figure}

As shown all the communication with master control unit is based on CAN bus. 
%!TEX root = ../Thesis.tex
\chapter{Motivation}
%!TEX root = ../Thesis.tex
\chapter{Implementation}
As previously mentioned the main control unit has be deployed on cRIO9033 - National Intruments' controller equipped with real-time processor and re-configurable FPGA unit. The reasoning for usage of this device was that it has more than sufficient computation power for the task and with additional modules provides support for basic CAN communication as well as high level CANOpen abstraction.


In my setup I have been using module NI9853 to directly receive and send a CAN messages and NI9881 for communication with the motor controllers based on CANOpen standard.

\section{Master control unit software implementation}
From the early beginning I have been developing the code of master control unit using event based programming paradigm. The reasoning of this is asynchronous, real-time nature of of such unit.

Most simple classification of subsystems can be based on communication protocol being used. 

\subsection{CAN integration}
To start with I had to implement a CAN library so it would be possible to make certain threads wait for message with predefined id, one from array of ids or just any message. 

To achieve it, I have a functionality deployed on FPGA (Field Programmable Gate Array) which simply translates input/output messages into the arrays of 8x unsigned integers. Read messages  are furthermore pushed into the FIFO register and interrupt on real-time processor is called.

Upon initiation of real-time programme additional process is spawned in the background which is responsible for communication with the FPGA. For sending a message it simply pushes data to FPGA and marks it as ready to send. When it comes to reading the process waits for interrupt from FPGA and acknowledges it after data read is finished. 

Then look up table is check for existence of conditional variable for given id, it is created if does not exist, and updated with new message. Additionally conditional variable containing id of last received message id is updated.

Based to this architecture I implemented functions which can wait for any message, one with certain id or one from ids set.

\subsection{Device operation states}
Although, during operation everything is event based the whole application itself can be considered a state machine where normal operation is just one of the states.

\todo{fig Init->boot-up seq->operational->exit}

\subsection{Operational state}
Implemented controller provides following functionality:
\begin{itemize}
    \item Log pedals, motor, pumps status and all occurred errors
    \item Updates motors target torque based on calculated electronic differential state (torque vectoring), power limitation and acceleration pedal position
    \item Control pumps power based on motors temperature and current target torque
    \item Subsystems connection loss procedures
    \item (Emergency) shutdown
    \item Read pedals, BMS, LV BMS and motors status.
\end{itemize}
\include{chapters/Testing}
%!TEX root = ../Thesis.tex
\chapter{Experimental results}
%!TEX root = ../Thesis.tex
\chapter{Future improvements}
\blinddocument
%!TEX root = ../Thesis.tex
\chapter{Conclusion}
Morbi pharetra ligula integer mollis mi nec neque ultrices vitae volutpat leo ullamcorper. In at tellus magna. Curabitur quis posuere purus. Cum sociis natoque penatibus et magnis dis parturient montes, nascetur ridiculus mus. Suspendisse tristique placerat feugiat. Aliquam vitae est at enim auctor ultrices eleifend a urna. Donec non tincidunt felis. Maecenas at suscipit orci.


\appendix
%!TEX root = ../Thesis.tex
\chapter{An Appendix}
Lorem ipsum dolor sit amet, consectetur adipisicing elit, sed do eiusmod tempor incididunt ut labore et dolore magna aliqua. Ut enim ad minim veniam, quis nostrud exercitation ullamco laboris nisi ut aliquip ex ea commodo consequat. Duis aute irure dolor in reprehenderit in voluptate velit esse cillum dolore eu fugiat nulla pariatur. Excepteur sint occaecat cupidatat non proident, sunt in culpa qui officia deserunt mollit anim id est laborum.
%!TEX root = ../Thesis.tex
%\chapter{Long chapter title with $\pi$ $π$ or π}
%\chapter{Long chapter title with \texorpdfstring{$\pi$ $π$ or π}{π π or π}}
\chapter{Long chapter \texorpdfstring{$\phi \land \sigma$}{phi and sigma} title with π, very long title, and also \texorpdfstring{$math = \sigma$}{math = sigma}}

My favorite RFC is \cite{rfc2549}.
Lorem ipsum dolor sit amet, consectetur adipisicing elit, sed do eiusmod tempor incididunt ut labore et dolore magna aliqua. Ut enim ad minim veniam, quis nostrud exercitation ullamco laboris nisi ut aliquip ex ea commodo consequat. Duis aute irure dolor in reprehenderit in voluptate velit esse cillum dolore eu fugiat nulla pariatur. Excepteur sint occaecat cupidatat non proident, sunt in culpa qui officia deserunt mollit anim id est laborum. \textsf{What about some inline sans serif text?}. $x = 4, y= 7$ which means that $\sqrt{4} = 2$.

Sans serif testing:
\begin{itemize}
    \item \textsf{$\pi$}
    \item \textsf{$π$}
    \item \textsf{π}
    \item \textsf{\emph{italic}}
    \item \textsf{\textbf{\emph{bold italic}}}
    \item \textsf{\textbf{bold}}
    \item \textsf{\texttt{teletype}}
    \item $\mathsf{Math\ Sans\ Serif}$
    \item \textsf{Text Sans Serif}
\end{itemize}



\backmatter
\printbibliography

\end{document}
